\section*{Highlights}
\begin{itemize}\itemsep1pt \parskip0pt
\item Re-photographing an unchanged street moves a VLM score two-thirds as far as changing it.
\item Six acquisition statistics predict almost none of that movement.
\item What is systematic in it is a tenth of a point, and it grows with the interval.
\item Camera geometry alone makes a VLM report change in 45\% of identical-scene pairs.
\item The usable unit is a few hundred paired observations, not a single sample point.
\end{itemize}

\begin{abstract}
Vision-language models are increasingly used to measure urban change from repeated street-level imagery, but their longitudinal reliability is not well understood. We test how much a perception score can change when the street itself does not undergo substantial redevelopment. Using 4,648 consecutive-epoch image pairs from 435 Google Street View standpoints across five US cities, we find that re-photographing the same street changes a perception score by 0.80 points on average, equivalent to 66.5\% of the difference between two different streets in the same city. Repeated model calls contribute almost no variation, while image re-encoding and prompt-order changes each account for about one fifth of the between-street difference. Six image statistics describing scattering, contrast, colour, exposure, sharpness and specularity explain almost none of the remaining epoch-to-epoch variation. A small systematic drift of about 0.1 points remains and increases with the interval between captures, consistent with minor physical changes not recorded by redevelopment labels. Controlled experiments further show that acquisition conditions can shift scores when camera and image properties are allowed to vary, and that the direction of these shifts depends on the model. In crowdsourced imagery, camera geometry alone causes a model to report physical change in 45\% of identical-scene pairs; normalising both images to a common virtual camera reduces this rate to 7.5\%. Despite poor reliability at the individual-location level, aggregation recovers a coherent redevelopment signal: changed streets are judged wealthier, better maintained, more enclosed and less green. These results show that vision-language measurement of urban change is reliable at the scale of hundreds of paired observations, but not at the scale of individual sample points.
\end{abstract}
